% theory_globalization_advanced.tex —— 潮流计算 通用理论：全局化与高级求解方法
% 本文件为理论层章节，遵循 docs/modules/THEORY_WRITING_GUIDE.md。

\section{全局化与高级求解方法理论}\label{sec:pf-th-glob}

\begin{theorynote}
本章为通用数学理论，按数值最优化与非线性方程组求解的标准文献体系撰写，覆盖：
线搜索（Armijo/Wolfe 条件、Zoutendijk 全局收敛定理）、信赖域子问题与柯西点、
Levenberg--Marquardt 与信赖域的等价关系、GLL 非单调线搜索、NCP/半光滑 Newton
理论（Fischer--Burmeister 函数、Clarke 广义 Jacobian、超线性收敛）与 PV/PQ 切换
的统一表述、HELM 全纯嵌入的复分析基础（幂级数、收敛半径、Padé 逼近与 Stahl
定理）、同伦延拓（概率一正则性、路径跟踪、转折点）、Newton--Krylov/JFNK
与非精确 Newton 的收敛理论，以及不动点迭代的 Anderson 加速（与 GMRES 的等价）；
并以平台求解器实测数据给出高级求解器收敛与可解性边界的数值基准。本章公式不要求逐式对应代码；与平台实现
（\srcpath{src/power\_flow/} 与 \srcpath{include/hacdcpf/power\_flow/}）的对应关系
见本章末"与实现的对应"表，超出实现的内容以"未实现扩展说明"框就地标记。
实现层的组件级转写见本手册"全局化与鲁棒求解器族"章。
\end{theorynote}

\subsection{记号与约定}\label{sec:pf-th-glob-notation}
\begin{itemize}
  \item $F:\mathbb R^n\to\mathbb R^n$：潮流失配向量（P/Q 平衡、直流功率平衡、
        换流器局部方程等的堆叠）；$x=(\theta,V_m,V_{dc},\ldots)$ 为堆叠状态；
        $J(x)=DF(x)$ 为 Jacobian；
  \item $\phi(x)=\tfrac12\lVert F(x)\rVert_2^2$：最小二乘 merit 函数，
        $g(x)=\nabla\phi(x)=J(x)^{\mathsf T}F(x)$；
  \item $d_k$ 或 $\Delta x_k$：第 $k$ 次迭代的搜索方向；$\alpha_k\in(0,1]$：步长；
        $\beta\in(0,1)$：回溯折减因子；$c_1,c_2$：Armijo/Wolfe 常数；
  \item $\Delta_k$：信赖域半径；$m_k(p)=\phi_k+g_k^{\mathsf T}p+\tfrac12
        p^{\mathsf T}B_k p$（$B_k=J_k^{\mathsf T}J_k$）：二次模型；
        $\rho_k=\mathrm{ared}_k/\mathrm{pred}_k$：增益比；
  \item $\lambda$ 在本章按上下文两用：信赖域/LM 小节中为 LM 阻尼参数，
        同伦小节中为同伦参数；$\mu\ge0$ 专用于 NCP 光滑化参数；
  \item $s\in\mathbb C$：HELM 嵌入参数；$V_i(s)$：母线 $i$ 电压的全纯嵌入；
        $[L/M]_f$：函数 $f$ 的 $(L,M)$ 阶 Padé 逼近；
  \item $\eta_k\in[0,1)$：非精确 Newton 强制项（forcing term）；
        $\mathcal K_m(A,v)$：$m$ 维 Krylov 子空间；
  \item $a\perp b$ 表示互补性 $a\ge0,\ b\ge0,\ ab=0$；
        $\partial F(x)$ 表示 Clarke 广义 Jacobian。
\end{itemize}

\subsection{线搜索与信赖域的全局化框架}\label{sec:pf-th-glob-framework}

潮流 Newton 法的局部二次收敛只在解的邻域内成立；全局化机制回答的问题是：
从任意初值出发，如何修正纯 Newton 步使迭代在全局意义下有收敛保证。
两类标准机制——线搜索与信赖域——都作用在 merit 函数 $\phi$ 上。

\begin{definition}[Armijo 充分下降条件]\label{sec:pf-th-glob-def-armijo}
给定下降方向 $d_k$（即 $g_k^{\mathsf T}d_k<0$）与常数 $c_1\in(0,1)$，
步长 $\alpha_k>0$ 满足 Armijo 条件，若
\begin{equation}
\phi(x_k+\alpha_k d_k)\le \phi(x_k)+c_1\alpha_k\,g_k^{\mathsf T}d_k .
\end{equation}
回溯线搜索从 $\alpha=1$ 出发按 $\alpha\leftarrow\beta\alpha$ 折减，
直到条件成立。由 Taylor 展开，只要 $d_k$ 是下降方向，充分小的 $\alpha$
必满足该条件，故回溯在有限步内终止。
\end{definition}

\begin{definition}[Wolfe 条件]\label{sec:pf-th-glob-def-wolfe}
步长 $\alpha_k$ 满足 Wolfe 条件，若 Armijo 条件成立且曲率条件
\begin{equation}
g(x_k+\alpha_k d_k)^{\mathsf T}d_k\ge c_2\,g_k^{\mathsf T}d_k,
\qquad 0<c_1<c_2<1
\end{equation}
成立。曲率条件排除"步长过小"的接受：它要求新点处方向导数的下降速率
已被足够多地"用掉"。
\end{definition}

\begin{theorem}[Zoutendijk 全局收敛定理]\label{sec:pf-th-glob-thm-zoutendijk}
设 $\phi$ 下有界，$g=\nabla\phi$ 在水平集
$\{x:\phi(x)\le\phi(x_0)\}$ 上 Lipschitz 连续（常数 $L$），
迭代 $x_{k+1}=x_k+\alpha_k d_k$ 的步长满足 Wolfe 条件。则
\begin{equation}
\sum_{k=0}^{\infty}\cos^2\theta_k\,\lVert g_k\rVert^2<\infty,
\qquad
\cos\theta_k=\frac{-g_k^{\mathsf T}d_k}{\lVert g_k\rVert\,\lVert d_k\rVert}.
\end{equation}
特别地，若方向满足\emph{夹角条件} $\cos\theta_k\ge\delta>0$，
则 $\lim_{k\to\infty}\lVert g_k\rVert=0$。
\end{theorem}

\begin{proof}
记 $s_k=\alpha_k d_k$。由曲率条件，
\[
(g_{k+1}-g_k)^{\mathsf T}s_k\ge(c_2-1)\,g_k^{\mathsf T}s_k .
\]
由 Lipschitz 连续性，
\[
(g_{k+1}-g_k)^{\mathsf T}s_k\le\lVert g_{k+1}-g_k\rVert\,\lVert s_k\rVert
\le L\,\alpha_k\lVert d_k\rVert\,\lVert s_k\rVert
=L\,\alpha_k\lVert d_k\rVert^2 .
\]
两式合并得
$\alpha_k\ge\dfrac{(c_2-1)g_k^{\mathsf T}d_k}{L\lVert d_k\rVert^2}
=\dfrac{(1-c_2)\cos\theta_k\lVert g_k\rVert}{L\lVert d_k\rVert}$，
其中用到 $-g_k^{\mathsf T}d_k=\cos\theta_k\lVert g_k\rVert\lVert d_k\rVert$。
代入 Armijo 条件：
\[
\phi_{k+1}\le\phi_k-c_1\alpha_k\,\cos\theta_k\lVert g_k\rVert\lVert d_k\rVert
\le\phi_k-\frac{c_1(1-c_2)}{L}\cos^2\theta_k\lVert g_k\rVert^2 .
\]
递推并对 $k$ 求和，由 $\phi$ 下有界，
$\frac{c_1(1-c_2)}{L}\sum_k\cos^2\theta_k\lVert g_k\rVert^2
\le\phi_0-\lim_k\phi_k<\infty$。
若 $\cos\theta_k\ge\delta>0$，则
$\delta^2\sum_k\lVert g_k\rVert^2<\infty$，故 $\lVert g_k\rVert\to0$。
\end{proof}

\begin{remark}[对非线性方程组的含义]\label{sec:pf-th-glob-rem-stationary}
Zoutendijk 定理保证的是 $g=J^{\mathsf T}F\to0$，即 $\phi$ 的驻点，
而不是 $F=0$。当 $J$ 亏秩时可能存在 $F\ne0$ 而 $J^{\mathsf T}F=0$ 的
\emph{非解驻点}（merit 函数的局部极小）。这是以最小二乘 merit 全局化
非线性方程组的固有局限；实践中以发散探测、重启动与回退链兜底。
Newton 方向 $d_k=-J_k^{-1}F_k$ 在 $J_k$ 非奇异且条件数有界时满足
夹角条件：$g_k^{\mathsf T}d_k=-F_k^{\mathsf T}F_k=-2\phi_k<0$。
\end{remark}

\begin{definition}[信赖域子问题与柯西点]\label{sec:pf-th-glob-def-tr}
信赖域方法在第 $k$ 步求解约束子问题
\begin{equation}
\min_{p\in\mathbb R^n}\ m_k(p)=\phi_k+g_k^{\mathsf T}p
+\tfrac12 p^{\mathsf T}B_k p,
\qquad \lVert p\rVert\le\Delta_k,
\end{equation}
$B_k=J_k^{\mathsf T}J_k$（Gauss--Newton 模型）。沿负梯度方向的约束极小点
称为\emph{柯西点}：
\begin{equation}
p_k^{C}=-\tau_k\frac{\Delta_k}{\lVert g_k\rVert}\,g_k,
\qquad
\tau_k=\begin{cases}1,&\lVert g_k\rVert^3\le\Delta_k\,g_k^{\mathsf T}B_k g_k
\ \text{或}\ g_k^{\mathsf T}B_k g_k\le0,\\[2pt]
\dfrac{\lVert g_k\rVert^3}{\Delta_k\,g_k^{\mathsf T}B_k g_k},&\text{其余}.\end{cases}
\end{equation}
即无约束最速下降极小点 $\bar\tau=\lVert g_k\rVert^3/(g_k^{\mathsf T}B_k g_k)$
对应的点未出界时取其本身，否则截断到边界。
\end{definition}

\begin{lemma}[柯西点的充分下降]\label{sec:pf-th-glob-lem-cauchy}
柯西点满足
\begin{equation}
m_k(p_k^{C})-m_k(0)\le-\tfrac12\lVert g_k\rVert\,
\min\Bigl(\Delta_k,\ \frac{\lVert g_k\rVert}{\lVert B_k\rVert}\Bigr).
\end{equation}
\end{lemma}

\begin{proof}
沿 $p(\tau)=-\tau g_k/\lVert g_k\rVert$（$\tau\in(0,\Delta_k]$）有
\[
m_k(p(\tau))-m_k(0)=-\tau\lVert g_k\rVert
+\frac{\tau^2}{2\lVert g_k\rVert^2}\,g_k^{\mathsf T}B_k g_k .
\]
\textbf{情形 1}：$g_k^{\mathsf T}B_k g_k\le0$。二次项非正，
$\tau_k=1$，故
$m_k-m_k(0)\le-\Delta_k\lVert g_k\rVert\le-\tfrac12\lVert g_k\rVert
\min(\Delta_k,\lVert g_k\rVert/\lVert B_k\rVert)$。

\textbf{情形 2}：$g_k^{\mathsf T}B_k g_k>0$ 且
$\bar\tau=\lVert g_k\rVert^3/(g_k^{\mathsf T}B_k g_k)\le\Delta_k$。
取 $\tau=\bar\tau$ 得
$m_k-m_k(0)=-\tfrac12\lVert g_k\rVert^4/(g_k^{\mathsf T}B_k g_k)
=-\tfrac12\lVert g_k\rVert\bar\tau$。
由 $g_k^{\mathsf T}B_k g_k\le\lVert B_k\rVert\lVert g_k\rVert^2$ 有
$\bar\tau\ge\lVert g_k\rVert/\lVert B_k\rVert$；结合
$\bar\tau\le\Delta_k$ 知
$\min(\Delta_k,\lVert g_k\rVert/\lVert B_k\rVert)\le\bar\tau$，结论成立。

\textbf{情形 3}：$g_k^{\mathsf T}B_k g_k>0$ 且 $\bar\tau>\Delta_k$。
取 $\tau=\Delta_k$，此时 $\Delta_k^2 g_k^{\mathsf T}B_k g_k/
(2\lVert g_k\rVert^2)<\tfrac12\Delta_k\lVert g_k\rVert$，故
$m_k-m_k(0)<-\tfrac12\Delta_k\lVert g_k\rVert
\le-\tfrac12\lVert g_k\rVert\min(\Delta_k,\lVert g_k\rVert/\lVert B_k\rVert)$。
\end{proof}

\begin{theorem}[信赖域全局收敛（陈述）]\label{sec:pf-th-glob-thm-tr}
设迭代只接受满足 $\mathrm{pred}_k=m_k(0)-m_k(p_k)$ 不小于柯西点下降之固定
比例（如 $\tfrac14$）的试探步，并按增益比
$\rho_k=\mathrm{ared}_k/\mathrm{pred}_k$ 调节半径
（$\rho_k$ 小则收缩并拒绝，$\rho_k$ 接近 1 则扩张）。在 $\phi$ 下有界、
$g$ Lipschitz、$\lVert B_k\rVert$ 一致有界的假设下，
$\liminf_{k\to\infty}\lVert g_k\rVert=0$；若再有 $\eta>0$ 的接受阈值
与聚点处 Jacobian 非奇异，则整个序列收敛且渐近取纯 Newton 步。
\end{theorem}

\begin{proof}[证明思路]
若 $\sum\Delta_k$ 在某一拒绝循环中保持有正下界以外的行为，则对足够小
半径，Taylor 余项使 $\rho_k\to1$，必出现接受步；于是无限次接受的子列上
$\mathrm{ared}_k\ge\eta\,\mathrm{pred}_k
\ge\frac{\eta}{4}\cdot\tfrac12\lVert g_k\rVert
\min(\Delta_k,\lVert g_k\rVert/\lVert B_k\rVert)$。
由 $\phi_k$ 单调下降且下有界，$\mathrm{ared}_k\to0$，从而
$\lVert g_k\rVert\min(\Delta_k,\cdot)\to0$；若 $\lVert g_k\rVert$ 有正下界
则 $\Delta_k\to0$，与"半径充分小必接受"矛盾。完整论证见
Nocedal \& Wright~\cite{pfglob-nw} 第 4 章。
\end{proof}

\begin{remark}[dogleg 折线步]\label{sec:pf-th-glob-rem-dogleg}
精确求解信赖域子问题代价高；Powell 的 dogleg
方法~\cite{pfglob-powell,pfglob-ds} 以折线近似解轨迹：当 Gauss--Newton 步
$p^{GN}=-(J^{\mathsf T}J)^{-1}g$ 在界内时取之；否则沿"柯西点 $\to$ Newton
步"的折线走到边界。折线包含柯西点，故引理~\ref{sec:pf-th-glob-lem-cauchy}
的下降估计成立，全局收敛性随之保证。
\end{remark}

\begin{theorem}[Levenberg--Marquardt 与信赖域的对应]\label{sec:pf-th-glob-thm-lmtr}
设 $\lambda>0$ 使 $B+\lambda I=J^{\mathsf T}J+\lambda I$ 正定，LM 步
\begin{equation}
p^{LM}(\lambda)=-\bigl(J^{\mathsf T}J+\lambda I\bigr)^{-1}J^{\mathsf T}F
\end{equation}
恰好是信赖域子问题在半径 $\Delta=\lVert p^{LM}(\lambda)\rVert$ 下的全局解；
反之，信赖域解在 $\lVert p^*\rVert=\Delta$（边界情形）时必具有 LM 形式。
\end{theorem}

\begin{proof}
信赖域子问题的最优性条件（Nocedal \&
Wright~\cite{pfglob-nw} 定理 4.1；Moré~\cite{pfglob-more}）：
$p^*$ 为全局解当且仅当存在 $\lambda^*\ge0$ 使
\begin{equation}
(B+\lambda^* I)p^*=-g,\qquad B+\lambda^* I\succeq0,\qquad
\lambda^*\bigl(\Delta-\lVert p^*\rVert\bigr)=0.
\end{equation}
对给定 $\lambda>0$，$B+\lambda I$ 正定，$p^{LM}(\lambda)$ 由第一式唯一
确定，互补条件以 $\Delta=\lVert p^{LM}(\lambda)\rVert$ 平凡成立，故
$p^{LM}(\lambda)$ 是该半径下的全局解。反向由互补条件：边界解
（$\lVert p^*\rVert=\Delta$）要求 $\lambda^*\ge0$ 且满足第一式，即
LM 形式。
\end{proof}

\begin{remark}\label{sec:pf-th-glob-rem-lmtr}
定理~\ref{sec:pf-th-glob-thm-lmtr} 表明 LM 阻尼与信赖域半径是同一族正则化
步长的两种参数化：$\lambda$ 单调地反比于有效半径
（$\lambda\to0$ 时 $p^{LM}\to p^{GN}$，$\lambda\to\infty$ 时步长趋于零且
方向趋于最速下降）。因此实现中以增益比 $\rho$ 调节 $\lambda$
（接受且 $\rho$ 大则 $\lambda/3$，拒绝则 $10\lambda$）与以 $\rho$ 调节
$\Delta$ 在理论上等价，差别只在参数化与数值细节。LM 的历史出处为
Levenberg~\cite{pfglob-levenberg} 与 Marquardt~\cite{pfglob-marquardt}。
\end{remark}

\begin{remark}[伪瞬态延拓（PTC）]\label{sec:pf-th-glob-rem-ptc}
PTC 将方程 $F(x)=0$ 视为伪动力系统 $\dot x=-F(x)$ 的稳态，以隐式 Euler
离散的 Newton 修正求解 $(J+\tfrac1{\delta t}I)\Delta x=-F$：
$\delta t\to0$ 时为强阻尼的显式 Euler 小步，$\delta t\to\infty$ 时退化为
纯 Newton 步。其收敛性分析（稳态吸引域内局部指数收敛到稳态）见
Kelley \& Keyes~\cite{pfglob-kelley-keyes}。SER（Self-Equalising Residual）
更新 $\delta t\leftarrow\delta t\,(r_{k-1}/r_k)^\gamma$ 是实践上有效的
启发式调度，其理论性质没有与 Zoutendijk 定理同等严格的结果。
\end{remark}

\begin{gapnote}
精确信赖域子问题求解（Moré--Sorensen 迭代）与大规模情形常用的
Steihaug--Toint 截断共轭梯度法在当前代码中均未实现；信赖域分支采用
dogleg 折线近似（经 \srcpath{include/hacdcpf/engine/kernel/globalization/globalization.hpp}
转发到 MIPSolvers 的 \srcpath{dogleg\_step}）。
此外实现只检查 Armijo 充分下降，未检查 Wolfe 曲率条件——对非线性方程组
merit 这是常见简化，但严格意义上
定理~\ref{sec:pf-th-glob-thm-zoutendijk} 的假设因此不被完整满足。
\end{gapnote}

\subsection{非单调线搜索（GLL）}\label{sec:pf-th-glob-gll}

单调 Armijo 搜索要求 merit 每步下降；当迭代序列需穿过 merit 面的
"窄谷"时，这会导致过小的步长与停滞。Grippo--Lampariello--Lucidi
（GLL）~\cite{pfglob-gll} 放松了单调性。

\begin{definition}[GLL 非单调 Armijo 条件]\label{sec:pf-th-glob-def-gll}
取窗口 $M\ge1$，定义非单调参考值
\begin{equation}
\phi_k^{R}=\max_{0\le j\le\min(k,M-1)}\phi_{k-j},
\end{equation}
接受条件为
\begin{equation}
\phi(x_k+\alpha_k d_k)\le\phi_k^{R}
+c_1\alpha_k\,g_k^{\mathsf T}d_k .
\end{equation}
$M=1$ 时退化为单调 Armijo 搜索；$M>1$ 时容许 merit 在窗口内暂时上升，
只要相对最近 $M$ 步的最坏值仍有充分下降。
\end{definition}

\begin{theorem}[GLL 全局收敛性（论证思路）]\label{sec:pf-th-glob-thm-gll}
在定理~\ref{sec:pf-th-glob-thm-zoutendijk} 的光滑性假设与夹角条件
$\cos\theta_k\ge\delta>0$ 下，采用 GLL 接受准则的迭代满足
$\liminf_k\lVert g_k\rVert=0$。
\end{theorem}

\begin{proof}[证明思路]
记参考下标 $\ell(k)$ 为窗口内取得最大值者：
$\phi_{\ell(k)}=\phi_k^{R}$。论证分三步（细节见
Grippo 等~\cite{pfglob-gll}）。
\textbf{(i)} 参考序列单调：由接受条件，
$\phi_{k+1}\le\phi_{\ell(k)}+c_1\alpha_k g_k^{\mathsf T}d_k
\le\phi_{\ell(k)}$，故下一段窗口的最大值不超过上一段，
$\{\phi_{\ell(k)}\}$ 非增且下有界，收敛到某 $\phi^R$。
\textbf{(ii)} 参考子列充分下降：沿参考下标构成的子列应用与
Zoutendijk 定理相同的估计，得
$\sum\cos^2\theta_{\ell(k)}\lVert g_{\ell(k)}\rVert^2<\infty$，
故 $\lVert g_{\ell(k)}\rVert\to0$。
\textbf{(iii)} 全序列由参考子列控制：每个 $k$ 距其参考下标不超过
$M$ 步；窗口内每步的回溯把 merit 增长控制在参考值的有限邻域内，
方向与步长的一致有界性再把 $g_k$ 与 $g_{\ell(k)}$ 的差控制在
$o(1)$，从而 $\liminf\lVert g_k\rVert=0$。
\end{proof}

\begin{remark}\label{sec:pf-th-glob-rem-gll}
窗口 $M$ 是鲁棒性与速度的工程折中：$M$ 大则更宽容、可能接受更多全
Newton 步（减少回溯次数与 Jacobian 重分解），但非单调震荡的幅度也更
大；$M$ 小则接近单调行为。潮流实践中 $M\in[3,10]$ 是常见取值。
\end{remark}

\subsection{Anderson 加速}\label{sec:pf-th-glob-anderson}

前推回代、FDPF 与 HELM 胚迭代等\emph{不动点}方法只具线性收敛
（第~\ref{sec:pf-th-cls}~章）。Anderson 加速~\cite{pfglob-anderson} 利用历史残差
外推，常能显著提升其收敛速度，且不需要 Jacobian。

\begin{definition}[Anderson 加速 AA($m$)]\label{def:pf-th-glob-anderson}
设不动点迭代 $x\mapsto g(x)$、残差 $f(x)=g(x)-x$。取记忆深度 $m$，第 $k$ 步令
$m_k=\min(m,k)$，用最近 $m_k{+}1$ 个残差解约束最小二乘
\begin{equation}
\alpha^{(k)}=\arg\min_{\sum_i\alpha_i=1}
\Bigl\lVert\sum_{i=0}^{m_k}\alpha_i\,f_{k-m_k+i}\Bigr\rVert_2,
\end{equation}
再取（阻尼因子 $\beta\in(0,1]$，$\beta{=}1$ 为无阻尼）
\begin{equation}
x_{k+1}=(1-\beta)\sum_{i=0}^{m_k}\alpha_i^{(k)}x_{k-m_k+i}
+\beta\sum_{i=0}^{m_k}\alpha_i^{(k)}g(x_{k-m_k+i}).
\end{equation}
$m=0$ 退化为（阻尼）Picard 迭代。约束最小二乘可等价化为对残差差分
$\Delta f_i=f_{i+1}-f_i$ 的无约束最小二乘，规模 $n\times m_k$，代价低。
\end{definition}

\begin{theorem}[线性情形与 GMRES 的等价（Walker--Ni）]\label{thm:pf-th-glob-anderson-gmres}
设 $g(x)=Ax+b$ 为仿射映射，不动点即解 $(I-A)x=b$。则无阻尼、全记忆（$m=\infty$）的
Anderson 加速在不停滞前提下，其第 $k$ 步迭代 $x_k$ 与 GMRES 施于 $(I-A)x=b$ 的第 $k$
步迭代重合，残差范数逐步相等 $\lVert f_k\rVert=\lVert r_k^{\mathrm{GMRES}}\rVert$
~\cite{pfglob-walkerni}。
\end{theorem}
\begin{proof}[证明要点]
无约束形式下 AA 更新为 $x_{k+1}=x_k+\beta f_k-(\Delta X_k+\beta\Delta F_k)\gamma^{(k)}$，
$\gamma^{(k)}=\arg\min_\gamma\lVert f_k-\Delta F_k\gamma\rVert$，其中
$\Delta X_k,\Delta F_k$ 收集历史差分 $\Delta x_i,\Delta f_i$。仿射映射使
$\Delta f_i=(A-I)\Delta x_i$，故 $\Delta F_k=(A-I)\Delta X_k$；最小二乘
$\min_\gamma\lVert f_k-\Delta F_k\gamma\rVert$ 是在 Krylov 仿射空间上极小化
$(I-A)$-残差，正是 GMRES 的最小残差性质（命题~\ref{sec:pf-th-glob-prop-gmres}）。
逐步对照生成子空间即得残差相等；完整对应见 Walker--Ni~\cite{pfglob-walkerni} 定理 2.2。
\end{proof}

\begin{remark}[非线性收敛与工程意义]\label{rem:pf-th-glob-anderson}
对压缩不动点 $g$（Lipschitz 常数 $\kappa<1$），AA($m$) 局部 r-线性收敛，渐近因子不劣于
Picard，实践中常显著更快~\cite{pfglob-tothkelley}。定理
\ref{thm:pf-th-glob-anderson-gmres} 揭示其本质：AA 是不动点迭代上的“类 GMRES”外推，
以 $O(m)$ 存储与一次小最小二乘换取超线性加速——对配网 BFS、FDPF 这类线性收敛且单步
廉价的方法尤为契合。$m$ 过大时 $\Delta F_k$ 病态，需列主元 QR 或 Tikhonov 正则与条件
数门限兜底。
\end{remark}

\begin{gapnote}
Anderson 加速在平台中\textbf{未实现}：不动点类求解器（\srcpath{run\_bfs\_sweep}、
\srcpath{FDPFSolver::solve}）按朴素 Picard/拟 Newton 迭代运行，HELM 用固定阶 Padé 而非
残差外推；平台的加速手段是 Newton 类（二次收敛）与 GMRES 内迭代
（\srcpath{src/power\_flow/newton\_krylov.cpp:gmres\_solve}），非不动点外推。
\end{gapnote}

\subsection{互补问题与半光滑 Newton 方法}\label{sec:pf-th-glob-ncp}

发电机无功限值（PV$\to$PQ 切换）与换流器限流都是"控制器在若干运行分支
间切换"的问题。统一它们的数学对象是互补问题。

\begin{definition}[互补问题与 NCP 函数]\label{sec:pf-th-glob-def-ncp}
标量互补条件写作 $0\le a\perp b\ge0$，即 $a\ge0,\ b\ge0,\ ab=0$。
函数 $\varphi:\mathbb R^2\to\mathbb R$ 称为 \emph{NCP 函数}，若
\begin{equation}
\varphi(a,b)=0\iff a\ge0,\ b\ge0,\ ab=0 .
\end{equation}
把逐分量互补条件堆叠进 $F(x)=0$，即用 NCP 函数把混合不等式-等式系统
化为单个非光滑方程组。
\end{definition}

\begin{proposition}[Fischer--Burmeister 函数]\label{sec:pf-th-glob-prop-fb}
Fischer--Burmeister（FB）函数~\cite{pfglob-fischer}
\begin{equation}
\Phi_{FB}(a,b)=\sqrt{a^2+b^2}-a-b
\end{equation}
是 NCP 函数。
\end{proposition}

\begin{proof}
$\sqrt{a^2+b^2}\ge\max(|a|,|b|)$。
（$\Leftarrow$）若 $a,b\ge0$ 且 $ab=0$，不妨 $b=0$，则
$\sqrt{a^2}-a-0=a-a=0$。
（$\Rightarrow$）设 $\Phi_{FB}(a,b)=0$，即 $\sqrt{a^2+b^2}=a+b$。
右端须非负，且 $a+b\ge\max(|a|,|b|)\ge|a|$ 给出 $b\ge0$，对称地
$a\ge0$。平方得 $a^2+b^2=a^2+2ab+b^2$，即 $ab=0$。
\end{proof}

\begin{proposition}[FB 的光滑性结构]\label{sec:pf-th-glob-prop-fbsmooth}
$\Phi_{FB}$ 在 $(0,0)$ 以外处处连续可微，
\begin{equation}
\nabla\Phi_{FB}(a,b)=\Bigl(\frac{a}{\sqrt{a^2+b^2}}-1,\
\frac{b}{\sqrt{a^2+b^2}}-1\Bigr);
\end{equation}
在原点不可微，但 $\Psi(a,b)=\tfrac12\Phi_{FB}(a,b)^2$ 在全平面连续
可微且 $\nabla\Psi(0,0)=0$（原点外 $\nabla\Psi=\Phi_{FB}\nabla\Phi_{FB}$，
其模被 $\sqrt2\,|\Phi_{FB}|$ 控制，随 $(a,b)\to0$ 趋于零）。
\end{proposition}

\begin{definition}[Clarke 广义 Jacobian 与半光滑性]\label{sec:pf-th-glob-def-clarke}
局部 Lipschitz 函数 $F:\mathbb R^n\to\mathbb R^m$ 在 $x$ 处的
B-次微分 $\partial_B F(x)$ 为所有"可微点列 Jacobian 的极限"之集合，
Clarke 广义 Jacobian 为其凸包
$\partial F(x)=\mathrm{co}\,\partial_B F(x)$。
$F$ 在 $x$ \emph{半光滑}，若对任意 $h\to0$ 与
$V\in\partial F(x+h)$ 有
\begin{equation}
F(x+h)-F(x)-Vh=o(\lVert h\rVert);
\end{equation}
若加强为 $O(\lVert h\rVert^2)$，则称\emph{强半光滑}。
分片光滑函数（如 $\min/\max/\Phi_{FB}$）皆强半光滑
（Facchinei \& Pang~\cite{pfglob-fp} 第 7 章）。
\end{definition}

\begin{example}[FB 在原点的广义 Jacobian]\label{sec:pf-th-glob-ex-fb}
$\Phi_{FB}$ 在原点的 B-次微分是单位圆周平移后的点集
$\{(\frac{u}{r}-1,\frac{v}{r}-1):u^2+v^2=r^2,\ r>0\}$ 的闭包
（以 $1/\sqrt2$ 为半径的圆），其 Clarke 广义 Jacobian 为对应圆盘。
实现对偶激活点 $(a,b)=(0,0)$ 选取对称元素
$(1/\sqrt2-1,\ 1/\sqrt2-1)$——即广义 Jacobian 的"中心"元，
避免对某一侧分支的偏向。
\end{example}

\begin{theorem}[半光滑 Newton 的局部超线性收敛（Qi--Sun）]\label{sec:pf-th-glob-thm-qisun}
设 $F$ 半光滑，$x^*$ 满足 $F(x^*)=0$，且 $\partial F(x^*)$ 中所有元素
非奇异（BD-正则性）。则广义 Newton 迭代
\begin{equation}
x_{k+1}=x_k-V_k^{-1}F(x_k),\qquad V_k\in\partial F(x_k)
\end{equation}
在 $x^*$ 的邻域内适定并 q-超线性收敛；若 $F$ 强半光滑，收敛是 q-二次的
（Qi \& Sun~\cite{pfglob-qisun}）。
\end{theorem}

\begin{proof}[证明思路]
BD-正则性与广义 Jacobian 的上半连续性给出 $x^*$ 邻域内
$\lVert V^{-1}\rVert$ 的一致界 $C$。于是
\[
x_{k+1}-x^*=V_k^{-1}\bigl[F(x^*)-F(x_k)-V_k(x^*-x_k)\bigr],
\]
括号内由半光滑性为 $o(\lVert x_k-x^*\rVert)$，故
$\lVert x_{k+1}-x^*\rVert=o(\lVert x_k-x^*\rVert)$，即超线性；
强半光滑时余项为 $O(\lVert x_k-x^*\rVert^2)$，得二次收敛。
邻域适定性由同一致界保证 $V_k$ 非奇异。
\end{proof}

\begin{remark}[BD-正则性不可由单点证书替代]\label{sec:pf-th-glob-rem-bd}
定理~\ref{sec:pf-th-glob-thm-qisun} 的假设要求解点邻域内
B-次微分的\emph{所有}元素非奇异，这比"所选广义 Jacobian 元非奇异"
强。运行中实测的速率商
$\lVert F_{k+1}\rVert/\lVert F_k\rVert$ 只是局部速率的经验证据，
不能反推 BD-正则性成立；退化分支交点（例如两个限值同时临界）处
所选广义 Jacobian 可能病态，实现以后向误差证书拒绝该步并回退全
稀疏 LU。
\end{remark}

\subsubsection{PV/PQ 切换的统一表述：中位数互补方程}

\begin{proposition}[中位数 NCP 方程的零点集]\label{sec:pf-th-glob-prop-median}
设母线电压偏差 $e=V_m-V_m^{set}$，无功裕度
$\underline q=Q_g-Q_g^{\min}$、$\bar q=Q_g-Q_g^{\max}$
（$\underline q>\bar q$）。方程
\begin{equation}
H(e,\underline q,\bar q)
=\max\bigl(\min(e,\underline q),\ \bar q\bigr)=0
\end{equation}
的零点集恰为以下三个物理分支的并：
\begin{equation}
\text{(i)}\ e=0,\ Q_g^{\min}\le Q_g\le Q_g^{\max};\qquad
\text{(ii)}\ Q_g=Q_g^{\min},\ e\ge0;\qquad
\text{(iii)}\ Q_g=Q_g^{\max},\ e\le0 .
\end{equation}
即"限值内为 PV（电压控制）"、触下/上限时转为 PQ 且电压向正确方向
松弛。
\end{proposition}

\begin{proof}
记 $u=\min(e,\underline q)$，则 $H=\max(u,\bar q)=0$ 当且仅当
$u\le0$，$\bar q\le0$，且两者至少其一为零。
\textbf{(i)} 若 $\bar q<0$（未触上限），须 $u=0$，即
$\min(e,\underline q)=0$：$e=0$ 且 $\underline q\ge0$，或
$\underline q=0$ 且 $e\ge0$（并入情形 (ii) 的边界）。
\textbf{(ii)} 若 $\underline q=0$（触下限），$u=\min(e,0)=0$
当且仅当 $e\ge0$，而 $\bar q=\underline q-(Q_g^{\max}-Q_g^{\min})<0$
自动满足。
\textbf{(iii)} 若 $\bar q=0$（触上限），须
$u=\min(e,\underline q)\le0$；此时 $\underline q>0$，
故须 $e\le0$。
三种情形之外：$\bar q<0$、$\underline q>0$、$e\ne0$ 时
$u=\min(e,\underline q)\ne0$ 且若 $e<0$ 则 $u<0$ 使 $H<0$，
$e>0$ 时 $u>0$ 使 $H>0$，均非零。
\end{proof}

\begin{remark}[与原始-对偶积极集的关系]\label{sec:pf-th-glob-rem-activeset}
把不等式约束的积极集显式枚举（PV$\to$PQ 批量切换加有界恢复审计）是
同一数学对象的另一算法化路径：对固定积极集解光滑系统、在收敛点处
按违反度更新积极集，其有限识别性质与半光滑 Newton 等价地建立在
互补结构上（Hintermüller--Ito--Kunisch~\cite{pfglob-hik}）。两者的
差别在工程侧：积极集切换改变方程布局与稀疏模式（需重做符号分析），
NCP 方程保持固定布局但引入非光滑点。
\end{remark}

\begin{definition}[CHKS 与 Kanzow 光滑化]\label{sec:pf-th-glob-def-smoothing}
CHKS 光滑化（Chen--Harker~\cite{pfglob-chenharker} 与
Kanzow~\cite{pfglob-kanzow} 的非内点延拓构造）以
\begin{equation}
\max\nolimits_\mu(a,b)=\tfrac12\bigl(a+b+\sqrt{(a-b)^2+4\mu^2}\bigr),
\qquad \mu>0
\end{equation}
替代 $\max$（$\min$ 对偶地），Kanzow 光滑 FB 为
\begin{equation}
\Phi_\mu(a,b)=\sqrt{a^2+b^2+2\mu^2}-a-b .
\end{equation}
两者在 $\mu>0$ 时光滑，$\mu\to0$ 时逐点（实际一致）收敛到对应非光滑
算子。
\end{definition}

\begin{proposition}[光滑化的一致逼近阶]\label{sec:pf-th-glob-prop-smootherr}
对所有 $(a,b)\in\mathbb R^2$ 与 $\mu>0$，
\begin{equation}
\bigl|\Phi_\mu(a,b)-\Phi_{FB}(a,b)\bigr|\le\sqrt2\,\mu .
\end{equation}
\end{proposition}

\begin{proof}
记 $r=\sqrt{a^2+b^2}$，则
$\Phi_\mu-\Phi_{FB}=\sqrt{r^2+2\mu^2}-r
=\dfrac{2\mu^2}{\sqrt{r^2+2\mu^2}+r}\le\dfrac{2\mu^2}{\sqrt2\,\mu}
=\sqrt2\,\mu$。
\end{proof}

\begin{remark}[光滑延拓的收敛逻辑]\label{sec:pf-th-glob-rem-smoothcont}
以 $\mu_k\downarrow0$ 的调度依次解光滑方程 $F_{\mu_k}(x)=0$、每个
$\mu$ 层以前一层解为初值，是标准的非内点延拓
（Kanzow~\cite{pfglob-kanzow}）：
命题~\ref{sec:pf-th-glob-prop-smootherr} 保证 $F_{\mu_k}\to F$ 一致，
若极限方程的解满足 BD-正则性，则足够小的 $\mu$ 层落在下一层的
Newton 吸引域内，整条延拓链以有限次逐层 Newton 求解完成。工程上
以两阶段调度（粗阶段慢退火、近收敛快退火）平衡鲁棒性与迭代数；
最终证书必须以 $\mu=0$ 的精确 NCP 残差重估，防止把光滑代理的收敛
误报为精确互补性。
\end{remark}

\subsection{全纯嵌入负荷流方法（HELM）的数学基础}\label{sec:pf-th-glob-helm}

HELM（Trias~\cite{pfglob-trias}）把潮流方程嵌入复参数 $s$ 的全纯族，
用解析延拓代替迭代。

\begin{definition}[全纯嵌入与幂级数]\label{sec:pf-th-glob-def-helm}
对 PQ 母线 $i$、PV 母线 $j$ 与平衡母线 $k$，嵌入取
\begin{align}
\sum_{\ell}Y_{i\ell}\,V_\ell(s)
&=s\,\frac{S_i^{*}}{V_i^{*}(s^{*})}-s\,Y_i^{sh}\,V_i(s),\\
\sum_{\ell}Y_{j\ell}\,V_\ell(s)
&=s\,\frac{P_j-\mathrm j Q_j(s)}{V_j^{*}(s^{*})},\qquad
V_j(s)V_j^{*}(s^{*})=W_j+s\bigl(|V_j^{set}|^2-W_j\bigr),\\
V_k(s)&=V_k^{spec},
\end{align}
其中 $W_j=|V_j(0)|^2$。$s=0$ 对应无注入\emph{胚}（germ）系统，
其解 $V(0)$ 由带平衡母线恒等行的齐次线性系统给出；$s=1$ 还原原
潮流方程。假设 $s\mapsto V(s)$ 在 $s=0$ 邻域全纯，展开
$V_i(s)=\sum_{n\ge0}c_i^{(n)}s^n$。
\end{definition}

\begin{proposition}[系数递推]\label{sec:pf-th-glob-prop-helmrec}
若胚系统矩阵非奇异，则各阶系数由一系列具有\emph{同一}系数矩阵的
线性系统递推确定：定义逆电压级数 $W_i(s)=1/V_i(s)=\sum_n
w_i^{(n)}s^n$，Cauchy 乘积 $V_i(s)W_i(s)=1$ 给出卷积递推
$\sum_{m=0}^{n}c_i^{(m)}w_i^{(n-m)}=\delta_{n0}$；
代入嵌入方程逐阶比较 $s^n$ 的系数，第 $n$ 阶右端只含
$c^{(0)},\ldots,c^{(n-1)}$（及 $w^{(n-1)}$）。
故全部系数经一次矩阵分解加多次回代求得。
\end{proposition}

\begin{proof}[证明思路]
嵌入方程除 $1/V_i(s^{*})$ 项外对 $V(s)$ 是线性的；$s=0$ 处
$V_i(0)\ne0$（胚非奇异）保证 $W_i(s)$ 同邻域全纯。代入幂级数并
比较系数是标准操作；线性项贡献第 $n$ 阶未知量
$c^{(n)}$ 的胚矩阵作用，含 $s$ 因子的项把注入右端推迟一阶，
故第 $n$ 阶系统闭式可解。严格性由隐函数定理的复形式保证
（胚处 Jacobian 即胚矩阵，非奇异）。
\end{proof}

\begin{theorem}[Cauchy--Hadamard 收敛半径]\label{sec:pf-th-glob-thm-ch}
幂级数 $\sum_n c^{(n)}s^n$ 的收敛半径
\begin{equation}
R=\Bigl(\limsup_{n\to\infty}\lVert c^{(n)}\rVert^{1/n}\Bigr)^{-1}
\end{equation}
等于 $s=0$ 到 $V(s)$ 最近奇点的距离；级数在 $|s|<R$ 内收敛到嵌入解，
在 $|s|>R$ 发散。
\end{theorem}

\begin{remark}[奇点的物理含义]\label{sec:pf-th-glob-rem-helmsing}
潮流方程的解支在鞍节点（saddle-node，电压失稳临界点）处两支合并，
映射 $s\mapsto V(s)$ 在该处出现平方根型分支点。因此 $R$ 度量的是
沿嵌入路径到最近鞍节点的"复距离"：重载系统分支点靠近 $s=1$，
$R<1$ 时原级数无法在 $s=1$ 处求值——这正是需要解析延拓的原因。
\end{remark}

\begin{definition}[Padé 逼近]\label{sec:pf-th-glob-def-pade}
函数 $f$ 的 $(L,M)$ 阶 Padé 逼近是满足
\begin{equation}
f(z)-\frac{P_L(z)}{Q_M(z)}=O(z^{L+M+1}),\qquad Q_M(0)=1
\end{equation}
的有理函数 $[L/M]_f=P_L/Q_M$（$\deg P_L\le L,\ \deg Q_M\le M$）。
其系数由幂级数系数的线性方程组（Toeplitz 结构）确定；
Padé 逼近的极点自动逼近 $f$ 的奇点，因而可在收敛圆外继续
$f$——这是有理逼近相对幂级数的本质优势
（Baker \& Graves-Morris~\cite{pfglob-baker}）。
\end{definition}

\begin{theorem}[Stahl 定理（陈述）]\label{sec:pf-th-glob-thm-stahl}
设 $f$ 在 $z=0$ 解析，在 $\mathbb C$ 中只有有限多个分支点型奇点。
则对角（及近对角）Padé 序列 $[n/n]_f$ 在 $\mathbb C$ 去掉一个
对数容量最小的割集系统后的紧子集上\emph{依容量收敛}到 $f$
（Stahl~\cite{pfglob-stahl}）；分母零点的极限分布集中在该极值割集上。
\end{theorem}

\begin{remark}[对 HELM 的操作性含义]\label{sec:pf-th-glob-rem-stahlops}
Stahl 定理把"Padé 是否可信"化为奇点结构问题：若嵌入解支在实轴段
$[0,1]$ 上无奇点，则（在有限分支点的假设下）近对角 Padé 在 $s=1$
处收敛到物理解；反之，若 $[0,1]$ 内存在鞍节点分支点，Padé 部分
极点被吸引到该处，逼近出现震荡/发散。实现据此把"Padé 逼近不收敛"
作为潮流不可行（或解支不可达）的实用判据。严格地说，依容量收敛
容许零容量例外集，且"发散必因不可行"的逆否口径依赖奇点结构假设，
故该判据是强证据而非纯数学证书；其实现层边界见本手册
"全局化与鲁棒求解器族"章的 HELM 小节。
\end{remark}

\begin{gapnote}
Stahl 定理的完整假设核验（嵌入映射奇点结构的辨识、极值割集计算）
与多解分支的系统性枚举在当前代码中均未实现；HELM 求解器采用固定的
对角 Padé 矩阵法与相邻逼近差的启发式收敛判据
（\srcpath{src/power\_flow/helm\_solver.cpp:pade\_at\_one}），
未做容量意义下的误差分析。
\end{gapnote}

\subsection{同伦延拓理论}\label{sec:pf-th-glob-homotopy}

\begin{definition}[同伦映射]\label{sec:pf-th-glob-def-homotopy}
设 $F(x)=0$ 为目标方程，$G(x)=0$ 为有已知解 $x_0$ 的平凡方程。
同伦是满足 $H(x,0)=G(x)$、$H(x,1)=F(x)$ 的连续映射
$H:\mathbb R^n\times[0,1]\to\mathbb R^n$。线性凸同伦
$H(x,\lambda)=(1-\lambda)G(x)+\lambda F(x)$ 是最简单的构造；
对潮流，等价做法是按 $\lambda$ 缩放负荷/注入：
$H(x,\lambda)=F_\lambda(x)$，其中 $\lambda=0$ 为零注入平启动系统，
$\lambda=1$ 为目标工况。
\end{definition}

\begin{theorem}[概率一正则性（参数化 Sard 定理，陈述）]\label{sec:pf-th-glob-thm-sard}
对带参数 $a\in\mathbb R^p$ 的光滑同伦 $H(x,\lambda;a)$，
若 $0$ 是 $(x,\lambda,a)$ 联合映射的正则值，则对 Lebesgue 几乎
所有 $a$，$0$ 是 $H(\cdot,\cdot;a)$ 的正则值，从而
$H^{-1}(0)$ 是 $\mathbb R^n\times[0,1]$ 中的一维光滑流形——
其连通分支或为闭曲线，或为两端落在 $\lambda\in\{0,1\}$ 截口上的
光滑路径；从 $(x_0,0)$ 出发的路径或以有限弧长到达 $\lambda=1$
截口（得到解），或发散至无穷。
（Allgower \& Georg~\cite{pfglob-allgower}。）
\end{theorem}

\begin{remark}[证明思想]\label{sec:pf-th-glob-rem-sard}
核心是 Sard 定理的参数化形式：临界值集测度为零，故"路径横截
$\lambda$ 截口"对几乎一切初值/参数成立。这把延拓法从"碰运气的
启发式"提升为有概率一保证的方法；剩下的失败模式只有路径无界
（目标问题无解或解支断裂）与数值跟踪丢失。
\end{remark}

\begin{definition}[预测-校正路径跟踪]\label{sec:pf-th-glob-def-pc}
增广系统 $H(x,\lambda)=0$ 的一维解流形以弧长 $s$ 参数化。
预测步沿切向 $(\dot x,\dot\lambda)$（由
$[D_xH\ \ D_\lambda H](\dot x,\dot\lambda)^{\mathsf T}=0$ 与归一化
条件确定）取 Euler 步；校正步在垂直于切向的超平面上以 Newton 法
回到路径（Newton 校正子的 Jacobian 为加边矩阵，非奇异性由路径
正则性保证）。
\end{definition}

\begin{remark}[转折点（turning points）]\label{sec:pf-th-glob-rem-turning}
路径可能相对 $\lambda$ 折叠：在转折点处 $D_xH$ 奇异、
$d\lambda/ds=0$，"固定 $\lambda$ 增量+逐点重解"的朴素策略在该处
失败（试探点无收敛解）。弧长参数化使增广 Jacobian 在普通转折点
仍非奇异，路径可平滑地"越过"折叠；鞍节点恰恰是潮流嵌入路径最
典型的转折点。按弧长连续编号的路径上，奇点指标（$D_xH$ 负
特征值计数）在简单转折点两侧差一，这也是 CPF 电压稳定分析的
理论入口（其工具化属本手册 CPF 章的主题）。
\end{remark}

\begin{remark}[步长自适应]\label{sec:pf-th-glob-rem-homstep}
实现采用的简化策略——成功则步长加倍（封顶）、失败则折半（低于
下界判失败）——是预测-校正框架去掉切向预测后的退化形式：以
上一收敛解直接作为下一 $\lambda$ 层的 Newton 初值（纯校正）。
在每层 Newton 收敛且层间距离小时等价于零阶预测。
\end{remark}

\begin{gapnote}
弧长参数化、加边 Newton 校正与转折点穿越在当前同伦实现中均未实现：
\srcpath{src/power\_flow/homotopy\_continuation.cpp:HomotopyContinuationSolver::solve}
按 $\lambda$ 单调递增逐层调用完整混合潮流求解
（\srcpath{detail::scale\_system\_by\_lambda} 缩放注入），
失败仅折半步长，路径折叠会被判为同伦失败而非穿越。
\end{gapnote}

\subsection{Newton--Krylov 与非精确 Newton 理论}\label{sec:pf-th-glob-nk}

大规模系统的 Newton 步 $J_k s_k=-F_k$ 不必精确求解；
Newton--Krylov 方法以 Krylov 子空间迭代近似求解（JFNK 中矩阵-向量
积 $Jv$ 可用差分 $[F(x+\epsilon v)-F(x)]/\epsilon$ 免 Jacobian 装配地
近似），其收敛理论由非精确 Newton 框架给出。

\begin{definition}[非精确 Newton 条件]\label{sec:pf-th-glob-def-inexact}
第 $k$ 步接受满足
\begin{equation}
\lVert F_k+J_k s_k\rVert\le\eta_k\lVert F_k\rVert,
\qquad \eta_k\in[0,1)
\end{equation}
的试探步 $s_k$，$\eta_k$ 称为强制项（forcing term）。
$\eta_k=0$ 退化为精确 Newton。
\end{definition}

\begin{theorem}[非精确 Newton 局部收敛（Dembo--Eisenstat--Steihaug）]\label{sec:pf-th-glob-thm-des}
设 $F$ 在解 $x^*$ 邻域连续可微，$J(x^*)$ 非奇异，$x_0$ 充分接近
$x^*$。则：
\begin{enumerate}
  \item 若 $\eta_k\le\bar\eta<1$，迭代局部 q-线性收敛；
  \item $x_k\to x^*$ q-超线性当且仅当
        $\lVert F_k+J_ks_k\rVert=o(\lVert F_k\rVert)$
       （特别地 $\eta_k\to0$ 是充分条件）；
  \item 若 $J$ Lipschitz 且
        $\lVert F_k+J_ks_k\rVert=O(\lVert F_k\rVert^2)$
        （如 $\eta_k=O(\lVert F_k\rVert)$），收敛为 q-二次。
\end{enumerate}
（Dembo, Eisenstat \& Steihaug~\cite{pfglob-des}；选择策略见
Eisenstat \& Walker~\cite{pfglob-ew}。）
\end{theorem}

\begin{proof}[证明思路]
记 $r_k=F_k+J_ks_k$ 为线性化残差，则
$x_{k+1}-x^*=x_k-x^*-J_k^{-1}(F_k-r_k)$。
对 $F(x^*)=0$ 作 Taylor 展开
$F_k=J_k(x_k-x^*)+O(\lVert x_k-x^*\rVert^2)$（$J$ Lipschitz），
得
\[
x_{k+1}-x^*
=J_k^{-1}r_k+O(\lVert x_k-x^*\rVert^2),
\]
故
$\lVert x_{k+1}-x^*\rVert
\le C\bigl(\eta_k\lVert F_k\rVert+\lVert x_k-x^*\rVert^2\bigr)$。
由 $\lVert F_k\rVert=O(\lVert x_k-x^*\rVert)$（$J(x^*)$ 非奇异
给出双边等价），三条结论依次成立。常数 $C$ 依赖
$\lVert J(x^*)^{-1}\rVert$ 与局部 Lipschitz 常数。
\end{proof}

\begin{remark}[强制项的自适应选择]\label{sec:pf-th-glob-rem-forcing}
远离解时取小 $\eta$ 是浪费（线性化模型本身误差大），近解时必须
$\eta\to0$ 以保超线性。Eisenstat--Walker~\cite{pfglob-ew} 给出两类
标准调度：$\eta_k=\lVert\phi_k-\phi_{k-1}^{\mathrm{pred}}\rVert/
\phi_{k-1}$ 型（反映模型保真度）与
$\eta_k=\gamma(\lVert F_k\rVert/\lVert F_{k-1}\rVert)^\alpha$ 型
（反映收敛速率），并配安全回退。定理~\ref{sec:pf-th-glob-thm-des}
只约束渐近行为，调度影响的是实际迭代数与内迭代总成本。
\end{remark}

\begin{proposition}[GMRES 的最小残差性质]\label{sec:pf-th-glob-prop-gmres}
GMRES~\cite{pfglob-saad-schultz} 第 $m$ 步在仿射空间
$x_0+\mathcal K_m(A,r_0)$ 中取使 $\lVert b-Ax\rVert_2$ 最小的
$x_m$。精确算术下 $m=n$ 必终止于真解；若 Arnoldi 过程在第 $j$ 步
满足 $h_{j+1,j}=0$（良性中断），则当前子空间已含精确解。
\end{proposition}

\begin{proof}
Arnoldi 关系 $AV_m=V_{m+1}\bar H_m$（$\bar H_m$ 为
$(m{+}1)\times m$ 上 Hessenberg 矩阵）与 $V_m$ 列正交性给出：对
$x=x_0+V_my$，
\[
\lVert b-Ax\rVert=\lVert r_0-AV_my\rVert
=\bigl\lVert\beta e_1-\bar H_m y\bigr\rVert,
\]
右端是小规模最小二乘问题，以 Givens 旋转在线更新，其解给出最小
残差。$m=n$ 时 $\mathcal K_n=\mathbb R^n$，最小残差为 0；
$h_{j+1,j}=0$ 时 $\mathcal K_j$ 是 $A$-不变子空间且含 $r_0$，
真解的残差校正落在其中。
\end{proof}

\begin{remark}[重启与预条件]\label{sec:pf-th-glob-rem-gmresrestart}
存储与正交化成本随 $m$ 线性增长，实用中取 GMRES($m$)：
每 $m$ 步以当前解重启。重启破坏有限终止性，可能停滞，故预条件
几乎总是必要：左预条件以 $P^{-1}A$ 替代 $A$ 构造 Krylov 子空间，
$P$ 应逼近 $A$ 且易解。对本平台的混合 AC/DC 系统，按实现头文件
契约注释，统一 Newton 的 Jacobian 呈块上三角结构
（AC 块 $A$、DC 块 $D$、耦合块 $B$、左下为零），此时"Schur 块
预条件子"实为精确块分解 $P=J$：预处理算子为恒等，GMRES 理论上
一步收敛；其价值不在 Krylov 加速，而在于分别分解 $A$、$D$ 以
降低填充与分解成本（Kelley~\cite{pfglob-kelley} 第 3、6 章的一般
讨论）。
\end{remark}

\begin{gapnote}
定理~\ref{sec:pf-th-glob-thm-des} 意义下的自适应强制项在当前实现中
未接线：GMRES 内迭代的相对残差容差取固定配置值
\fld{gmres\_tol}（\srcpath{src/power\_flow/newton\_krylov.cpp:gmres\_solve}），
NK 步的接受另加宽松阈值与有限性检查
（\srcpath{src/power\_flow/newton\_krylov.cpp:newton\_krylov\_step}），
因此外层只享有结论 (i) 的线性收敛口径；免 Jacobian 的差分矩阵-向量
积（严格 JFNK）亦未采用，矩阵-向量积直接作用于已装配的 Jacobian。
\end{gapnote}

\subsection{数值算例与基准}\label{sec:pf-th-glob-numexp}

本节数字由 \srcpath{tools/pf\_doc\_benchmark.cpp}（模式 \texttt{robust}/\texttt{stress}）
在 MATPOWER 算例上经平台公共 API 实测（平启动，Apple~M4）。

\subsubsection{高级求解器在标准算例上的收敛}
表~\ref{tab:pf-th-glob-robust} 给出平启动下\emph{关闭}全局化的纯 Newton、同伦延拓、HELM
与 Newton--Krylov 的最终残差。四者在 $9\sim57$ 母线算例上均达到 $10^{-8}$ 容差；
\textbf{HELM 在 case118/case300 上止步于 $\sim10^{-7}$}（固定对角 Padé 阶的精度上限，未达
$10^{-8}$，表中标 $\dagger$），与注记~\ref{sec:pf-th-glob-rem-stahlops} 及实现层 HELM
边界一致。纯 Newton 在标准算例平启动下即收敛，说明全局化阶梯的价值在于\emph{更难}的
工况（见下）而非常规潮流。
\begin{table}[H]\centering\footnotesize
\caption{高级求解器在标准算例上的最终残差（平启动，$\mathrm{tol}=10^{-8}$）}
\label{tab:pf-th-glob-robust}
\begin{tabular}{@{}lrllll@{}}
\toprule
算例 & 母线数 & 纯 Newton & 同伦延拓 & HELM & Newton--Krylov\\
\midrule
case9   & 9   & $1.7\times10^{-14}$ & $3.4\times10^{-9}$  & $8.2\times10^{-10}$      & $1.7\times10^{-14}$\\
case14  & 14  & $1.3\times10^{-10}$ & $5.8\times10^{-11}$ & $1.6\times10^{-10}$      & $1.3\times10^{-10}$\\
case30  & 30  & $9.6\times10^{-10}$ & $1.4\times10^{-12}$ & $3.5\times10^{-11}$      & $9.6\times10^{-10}$\\
case57  & 57  & $3.4\times10^{-12}$ & $9.9\times10^{-9}$  & $1.6\times10^{-10}$      & $3.4\times10^{-12}$\\
case118 & 118 & $8.5\times10^{-14}$ & $8.5\times10^{-14}$ & $8.9\times10^{-8}\,\dagger$ & $8.5\times10^{-14}$\\
case300 & 300 & $4.3\times10^{-13}$ & $6.2\times10^{-13}$ & $2.5\times10^{-7}\,\dagger$ & $4.3\times10^{-13}$\\
\bottomrule
\end{tabular}
\end{table}

\subsubsection{负荷加载与可解性边界}
表~\ref{tab:pf-th-glob-stress} 把算例负荷按因子 $\gamma$ 整体放大（发电不变、slack 吸收），
记录纯 Newton、带鲁棒升级的 Newton 与同伦三者的收敛性。三者在每个 $\gamma$ 处\emph{一致}
收敛或一致失败（故合并为一列）：case9 的边界落在 $\gamma\in(2.2,2.5)$、case118 在
$(1.5,1.8)$。这印证：越过该边界的失败是\textbf{物理无解}（鼻点之外），而非求解器鲁棒性
不足——对应同伦理论中路径“发散至无穷”的失败模式（定理~\ref{sec:pf-th-glob-thm-sard}）
与可解性充分条件越界（定理~\ref{thm:pf-th-fund-solvability}）。全局化阶梯能扩大\emph{可
达}解集，却不能制造不存在的解。
\begin{table}[H]\centering\footnotesize
\caption{负荷加载因子 $\gamma$ 下的收敛性（纯 Newton、鲁棒 Newton、同伦三者一致）}
\label{tab:pf-th-glob-stress}
\begin{tabular}{@{}lccccccccc@{}}
\toprule
$\gamma$ & 1.0 & 1.2 & 1.5 & 1.8 & 2.0 & 2.2 & 2.5 & 2.8 & 3.0\\
\midrule
case9   & Y & Y & Y & Y & Y & Y & N & N & N\\
case118 & Y & Y & Y & N & N & N & N & N & N\\
\bottomrule
\end{tabular}
\end{table}

\subsection{与实现的对应}\label{sec:pf-th-glob-map}
\begin{center}\footnotesize
\begin{longtable}{@{}L{6.8cm}L{8.6cm}@{}}
\toprule
理论条目 & 实现状态\\
\midrule
\endfirsthead
\toprule
理论条目 & 实现状态\\
\midrule
\endhead
Armijo 充分下降 + 回溯 &
\implfull{include/hacdcpf/power\_flow/globalization/nonmonotone\_linesearch.hpp:NonmonotoneLineSearch::search}（$M{=}1$ 时即单调 Armijo；主循环挂接见 src/power\_flow/newton\_solver.cpp:NewtonSolver::solve）\\
Wolfe 曲率条件 & \implnone\\
Zoutendijk 全局收敛（定理~\ref{sec:pf-th-glob-thm-zoutendijk} 的渐近保证） &
\implpart{实现以有限回溯上限 + 停滞检测 + 五级失败回退链兜底（RobustNonlinearOptions 停滞与回退字段），无渐近收敛性证书}\\
信赖域子问题、柯西点与 dogleg（定义~\ref{sec:pf-th-glob-def-tr}、
引理~\ref{sec:pf-th-glob-lem-cauchy}） &
\implfull{include/hacdcpf/engine/kernel/globalization/globalization.hpp:dogleg\_step}（转发到 MIPSolvers；增益比半径调节同文件 trust\_region\_update，调用侧为 newton\_solver.cpp 的 TrustRegion 分支）\\
精确信赖域子问题（Moré--Sorensen）、Steihaug--Toint CG & \implnone\\
LM 正规方程 $(J^{\mathsf T}J+\lambda I)\Delta x=-J^{\mathsf T}F$
与增益比调 $\lambda$（定理~\ref{sec:pf-th-glob-thm-lmtr} 的参数化） &
\implpart{LM 恢复步内联实现于 src/power\_flow/newton\_solver.cpp:NewtonSolver::solve（线搜索耗尽后至多 4 次 $\lambda\times10$ 尝试）；头文件 LMController::update 是预留控制器，未被主循环消费}\\
PTC 与 SER 伪时间步（注记~\ref{sec:pf-th-glob-rem-ptc}） &
\implfull{include/hacdcpf/power\_flow/globalization/lm\_trust\_region.hpp:PtcSerController::update}（PTC 分支装配 $J+I/\delta t$，挂接 NewtonSolver::solve）\\
GLL 非单调窗口参考值（定义~\ref{sec:pf-th-glob-def-gll}） &
\implfull{include/hacdcpf/power\_flow/globalization/nonmonotone\_linesearch.hpp:NonmonotoneLineSearch}（push\_merit/reference\_merit，窗口默认 $M=5$）\\
FB 函数与广义 Jacobian 元素
（命题~\ref{sec:pf-th-glob-prop-fb}、例~\ref{sec:pf-th-glob-ex-fb}） &
\implfull{include/hacdcpf/power\_flow/ncp\_functions.hpp:fischer\_burmeister}（原点对称元见同文件 fischer\_burmeister\_jacobian）\\
中位数 NCP 统一 PV/PQ（命题~\ref{sec:pf-th-glob-prop-median}） &
\implfull{include/hacdcpf/power\_flow/ncp\_functions.hpp:pv\_pq\_ncp}（广义导数 pv\_pq\_ncp\_jacobian，并列处取相邻单侧斜率平均）\\
CHKS/Kanzow 光滑化与 $\mu$ 延拓（定义~\ref{sec:pf-th-glob-def-smoothing}） &
\implfull{include/hacdcpf/power\_flow/ncp\_functions.hpp:smooth\_fb}（pv\_pq\_smooth\_ncp；两阶段退火调度在 NewtonSolver::solve 内，参数 RobustNonlinearOptions::enable\_smooth\_ncp 系列）\\
半光滑 Newton 超线性/二次收敛（定理~\ref{sec:pf-th-glob-thm-qisun}） &
\implpart{求解器按广义 Newton 步运行并上报实测速率商；BD-正则性不做事先验证，病态所选块以后向误差证书拒绝并回退全稀疏 LU（见 vsc\_limit\_ncp 契约）}\\
VSC 限流的互补表述（同一 NCP 框架的设备级实例） &
\implfull{src/power\_flow/vsc\_limit\_ncp.cpp:evaluate\_vsc\_limit\_ncp}（median\_ncp；限流圆盘与优先级分支见 vsc\_limit\_command）\\
HELM 嵌入与系数递推（定义~\ref{sec:pf-th-glob-def-helm}、
命题~\ref{sec:pf-th-glob-prop-helmrec}） &
\implfull{src/power\_flow/helm\_solver.cpp:HelmSolver::solve}（胚矩阵 build\_germ\_matrix、$|V|^2$ 约束系数 magnitude\_squared\_coefficient）\\
Padé 解析延拓（定义~\ref{sec:pf-th-glob-def-pade}） &
\implfull{src/power\_flow/helm\_solver.cpp:pade\_at\_one}（对角 Padé 矩阵法在 $s=1$ 求值）\\
Stahl 定理层面的收敛性分析（定理~\ref{sec:pf-th-glob-thm-stahl}） &
\implpart{实现以相邻 Padé 逼近差为启发式收敛判据、以发散作不可行性实用判据；无容量/割集层面的理论核验}\\
同伦映射与注入缩放（定义~\ref{sec:pf-th-glob-def-homotopy}） &
\implfull{src/power\_flow/homotopy\_continuation.cpp:detail::scale\_system\_by\_lambda}（逐层驱动 HomotopyContinuationSolver::solve）\\
概率一正则性与预测-校正跟踪
（定理~\ref{sec:pf-th-glob-thm-sard}、定义~\ref{sec:pf-th-glob-def-pc}） &
\implpart{实现为零阶预测（上层解直接作初值）+ 逐层完整重解 + 步长成败折半/加倍（HomotopyState）；无正则性检验}\\
弧长参数化与转折点穿越（注记~\ref{sec:pf-th-glob-rem-turning}） & \implnone\\
非精确 Newton 条件与 DES 收敛理论
（定义~\ref{sec:pf-th-glob-def-inexact}、
定理~\ref{sec:pf-th-glob-thm-des}） &
\implpart{内层残差容差为固定 \fld{gmres\_tol} 而非自适应强制项（newton\_krylov.cpp:newton\_krylov\_step），外层仅享线性收敛口径}\\
GMRES($m$)：MGS Arnoldi + Givens 最小二乘
（命题~\ref{sec:pf-th-glob-prop-gmres}） &
\implfull{src/power\_flow/newton\_krylov.cpp:gmres\_solve}（重启、真残差复核、良性中断处理齐全）\\
Schur 块预条件（注记~\ref{sec:pf-th-glob-rem-gmresrestart}） &
\implfull{src/power\_flow/newton\_krylov.cpp:SchurBlockPreconditioner::build}（apply 实现两段块回代；块结构契约见 newton\_krylov.hpp 类注释）\\
Anderson 加速与 GMRES 等价（定义~\ref{def:pf-th-glob-anderson}、
定理~\ref{thm:pf-th-glob-anderson-gmres}） &
\implnone（不动点求解器用朴素 Picard/拟 Newton，无残差外推；加速手段为 Newton/GMRES）\\
高级求解器与可解性边界数值基准（表~\ref{tab:pf-th-glob-robust}、
\ref{tab:pf-th-glob-stress}） &
\implfull{tools/pf\_doc\_benchmark.cpp}（模式 robust/stress）\\
\bottomrule
\end{longtable}
\end{center}

\subsection{参考文献}
{\renewcommand{\section}[2]{}%
\begin{thebibliography}{99}\small
\bibitem{pfglob-nw} J.~Nocedal and S.~J. Wright, \emph{Numerical
Optimization}, 2nd ed., Springer, 2006.
\bibitem{pfglob-ds} J.~E. Dennis and R.~B. Schnabel, \emph{Numerical
Methods for Unconstrained Optimization and Nonlinear Equations}, SIAM,
1996.
\bibitem{pfglob-powell} M.~J.~D. Powell, A hybrid method for nonlinear
equations, in \emph{Numerical Methods for Nonlinear Algebraic Equations}
(P.~Rabinowitz, ed.), Gordon and Breach, 1970.
\bibitem{pfglob-levenberg} K.~Levenberg, A method for the solution of
certain non-linear problems in least squares, \emph{Quarterly of Applied
Mathematics}, 2, 1944.
\bibitem{pfglob-marquardt} D.~W. Marquardt, An algorithm for least-squares
estimation of nonlinear parameters, \emph{SIAM Journal on Applied
Mathematics}, 11(2), 1963.
\bibitem{pfglob-more} J.~J. Mor\'e, The Levenberg--Marquardt algorithm:
implementation and theory, in \emph{Numerical Analysis}, Lecture Notes in
Mathematics 630, Springer, 1978.
\bibitem{pfglob-gll} L.~Grippo, F.~Lampariello and S.~Lucidi, A nonmonotone
line search technique for Newton's method, \emph{SIAM Journal on Numerical
Analysis}, 23(4), 1986.
\bibitem{pfglob-fp} F.~Facchinei and J.-S. Pang, \emph{Finite-Dimensional
Variational Inequalities and Complementarity Problems}, Vol.~I, Springer,
2003.
\bibitem{pfglob-fischer} A.~Fischer, A special Newton-type optimization
method, \emph{Optimization}, 24, 1992.
\bibitem{pfglob-qisun} L.~Qi and J.~Sun, A nonsmooth version of Newton's
method, \emph{Mathematical Programming}, 58, 1993.
\bibitem{pfglob-kanzow} C.~Kanzow, Some noninterior continuation methods
for linear complementarity problems, \emph{SIAM Journal on Matrix Analysis
and Applications}, 17(4), 1996.
\bibitem{pfglob-chenharker} B.~Chen and P.~T. Harker, A non-interior-point
continuation method for linear complementarity problems, \emph{SIAM Journal
on Matrix Analysis and Applications}, 14(4), 1993.
\bibitem{pfglob-hik} M.~Hinterm\"uller, K.~Ito and K.~Kunisch, The
primal-dual active set strategy as a semismooth Newton method,
\emph{SIAM Journal on Optimization}, 13(3), 2002.
\bibitem{pfglob-trias} A.~Trias, The holomorphic embedding load flow
method, in \emph{Proceedings of the IEEE PES General Meeting}, San Diego,
2012.
\bibitem{pfglob-stahl} H.~Stahl, On the convergence of generalized Pad\'e
approximants, \emph{Constructive Approximation}, 5, 1989.
\bibitem{pfglob-baker} G.~A. Baker and P.~Graves-Morris, \emph{Pad\'e
Approximants}, 2nd ed., Cambridge University Press, 1996.
\bibitem{pfglob-allgower} E.~L. Allgower and K.~Georg, \emph{Numerical
Continuation Methods: An Introduction}, Springer, 1990.
\bibitem{pfglob-des} R.~S. Dembo, S.~C. Eisenstat and T.~Steihaug, Inexact
Newton methods, \emph{SIAM Journal on Numerical Analysis}, 19(2), 1982.
\bibitem{pfglob-ew} S.~C. Eisenstat and H.~F. Walker, Choosing the forcing
terms in an inexact Newton method, \emph{SIAM Journal on Scientific
Computing}, 17(1), 1996.
\bibitem{pfglob-saad-schultz} Y.~Saad and M.~H. Schultz, GMRES: a
generalized minimal residual algorithm for solving nonsymmetric linear
systems, \emph{SIAM Journal on Scientific and Statistical Computing},
7(3), 1986.
\bibitem{pfglob-kelley} C.~T. Kelley, \emph{Iterative Methods for Linear
and Nonlinear Equations}, SIAM, 1995.
\bibitem{pfglob-kelley-keyes} C.~T. Kelley and D.~E. Keyes, Convergence
analysis of pseudo-transient continuation, \emph{SIAM Journal on Numerical
Analysis}, 35(2), 1998.
\bibitem{pfglob-anderson} D.~G. Anderson, Iterative procedures for nonlinear
integral equations, \emph{Journal of the ACM}, 12(4), 1965.
\bibitem{pfglob-walkerni} H.~F. Walker and P.~Ni, Anderson acceleration for
fixed-point iterations, \emph{SIAM Journal on Numerical Analysis}, 49(4), 2011.
\bibitem{pfglob-tothkelley} A.~Toth and C.~T. Kelley, Convergence analysis for
Anderson acceleration, \emph{SIAM Journal on Numerical Analysis}, 53(2), 2015.
\end{thebibliography}}
